% guide.tex - the Ultimate SDK Programmer's Reference Guide.
%
% Build:  make -C guide        (or: latexmk -pdf guide.tex)
%
% The chapters are separate files under chapters/, and the order here is the
% order of the book. Front matter first, then getting started, then one chapter
% per language, then the services, then the appendices.
%
% SPDX-License-Identifier: MIT

\documentclass[10pt,openany]{book}

% A small page, close to the original's trim. A narrow measure is most of why
% the Commodore guides are easy to read, and it is worth keeping even though a
% programmer's book carries wider code than a user's guide does.
\usepackage[papersize={6.0in,9.0in},
            top=0.85in,bottom=0.85in,inner=0.80in,outer=0.65in]{geometry}

\usepackage{ultimatesdk}
\usepackage[hidelinks,pdfusetitle]{hyperref}

\title{Ultimate SDK Programmer's Reference Guide}
\author{}
\date{}

\begin{document}

% ---------------------------------------------------------------------------
% Title page: the blue cover of the original, set plainly.
% ---------------------------------------------------------------------------
\frontmatter
\thispagestyle{empty}
\AddToShipoutPictureBG*{\AtPageLowerLeft{%
  \color{guideblue}\rule{\paperwidth}{\paperheight}}}
\null\vspace*{0.20\paperheight}
{\color{white}\headingfont\centering
 {\fontsize{30}{34}\selectfont ULTIMATE SDK}\\[0.9\baselineskip]
 {\fontsize{20}{24}\selectfont PROGRAMMER'S}\\[0.25\baselineskip]
 {\fontsize{20}{24}\selectfont REFERENCE GUIDE}\par}
\vfill
{\color{white}\centering\normalfont\small
 For the 1541 Ultimate-II, Ultimate 64, Ultimate 64 Elite\\
 and the Commodore 64 Ultimate\par}
\vspace{0.10\paperheight}
\clearpage

\thispagestyle{empty}
\null\vfill
\noindent\small
First edition.\\[0.4\baselineskip]
This guide documents the Ultimate SDK, a software development kit for the
Commodore~64 and the Ultimate family of cartridges and machines.\\[0.6\baselineskip]
Every figure, timing and behaviour in this book was measured against real
hardware running firmware 3.15 unless it says otherwise. Where the firmware
does something other than its documentation describes, this guide follows the
firmware and says so.\\[0.6\baselineskip]
Released under the MIT licence. Commodore, the Commodore logo and Commodore~64
are trademarks of their respective owners; this guide is not affiliated with or
endorsed by them.
\clearpage

\tableofcontents

% ---------------------------------------------------------------------------
\cleardoublepage
\section*{Introduction}
\addcontentsline{toc}{section}{Introduction}

Welcome. If you have an Ultimate cartridge or an Ultimate~64, you own a
Commodore~64 that can reach a great deal more than a Commodore~64 ever could
--- files on a USB stick, sixteen megabytes of RAM expansion, a processor that
runs up to forty-eight times faster than the one Commodore shipped, and a
network connection to the rest of the world. This guide is about writing
programs that use all of it.

The \sdk{} is a library, and it comes in the language you already write in.
If you write BASIC, it arrives as new keywords you type at the \bas{ready.}
prompt. If you write C with cc65, it is a header and a library to link. If you
write assembly, it is a set of entry points and a documented calling
convention. And if you use a toolchain that cannot link any of those, there is
a standalone binary with a jump table that any assembler can call.

The four routes are not four implementations. They are one implementation,
reached four ways, and every one of them is tested against the same hardware.

\subsection*{What is in it}

The SDK started as a wrapper around the Ultimate Command Interface --- the
hardware mailbox at \cmd{\$DF1C} that a program uses to ask the cartridge to do
something --- and it has grown past that. Three of the things it does are not
command-interface operations at all, because the machine offers them a
different way:

\begin{itemize}[leftmargin=1.4em,itemsep=0.25em,topsep=0.4em]
  \item \textbf{Files and directories}, over the command interface.
  \item \textbf{Loading and saving}, taking the firmware's fast path when
        there is one.
  \item \textbf{The RAM expansion}, both by direct memory access --- which is
        not a command at all, it is the REU's own registers --- and file to
        expansion without the bytes passing through the C64.
  \item \textbf{Processor speed} on an Ultimate~64, which is memory-mapped I/O
        and has no command behind it either.
  \item \textbf{The screen palette}, all sixteen colours, live.
  \item \textbf{Network sockets and HTTP}, so a C64 can fetch a file from a
        web server.
\end{itemize}

\subsection*{How to read this}

Chapter~1 gets you running something in the next ten minutes, whichever
language you write in. Chapters~2 to~5 take one language each --- BASIC, C,
assembly, and everything else --- and are meant to be read start to finish once
and then dipped into. The chapters after that take one service each and can be
read in any order, or not at all until you need them.

Everything is small enough to type in. The programs in this book are short on
purpose: a listing you can enter in a minute and watch work teaches more than
one you copy from a disk and take on trust.

\begin{note}
This guide describes behaviour that was measured, not behaviour that was
assumed. Where a figure appears --- how many cycles something costs, how long a
connection takes to fail --- it came off a real machine, and the guide says
which one.
\end{note}

\subsection*{A word about the name}

The library is called the Ultimate SDK rather than the Ultimate~64 SDK because
most of it works on every Ultimate from the 1541~Ultimate-II upward, and not
only on the Ultimate~64. Where something needs a particular machine --- the
turbo registers need an Ultimate~64 --- the guide says so at the point it
matters, and the SDK itself tells you at run time rather than guessing.

\mainmatter

% Chapter 1 - Getting Started.

\guidechapter{GETTING STARTED}{%
  \item What the SDK requires
  \item Enabling the command interface
  \item Choosing and checking a binding
  \item Recovering from common failures
}

\section{What The SDK Requires}

The SDK runs on a C64 with Ultimate hardware whose command interface is mapped
at \cmd{\$DF1B--\$DF1F}. It is supported on the 1541~Ultimate-II and later
products.

There is no blanket minimum firmware version. Targets were added over time, so
programs call \cmd{ultimate\_detect()} and test the capability they need. The
palette is newer than firmware 3.15; HTTP appears in 3.15.

The control target predates the palette commands, so its presence is not a
palette capability test. Call a palette function and handle
\cmd{ULTIMATE\_ERR\_NOT\_SUPPORTED}; the SDK returns that result when older
firmware answers that it does not know the command.

\evidence{\cmd{docs/compatibility.md}; \cmd{docs/uci.md};
\cmd{include/ultimate.h}}

\section{Enable The Command Interface}

Open the Ultimate configuration menu and set \bas{command interface} to
\bas{enabled}. The exact key used to open that menu depends on the product and
its configuration, so the SDK does not prescribe one.

With the interface disabled, the five registers are not mapped and
\cmd{ultimate\_init()} returns \cmd{ULTIMATE\_ERR\_NO\_DEVICE}.

REU transfers and Ultimate~64 turbo control use their own registers, not the
command interface. Each still depends on its corresponding machine setting.

\evidence{\cmd{docs/compatibility.md};
\cmd{tests/hardware/hwtest.py}; \cmd{tools/test\_registers.py}}

\section{Choose A Binding}

\begin{reftable}{lp{0.63\linewidth}}
BASIC V2 & \cmd{make wedge}; load \cmd{src/basic/uci.prg}, or mount
  \cmd{src/basic/uci.crt} \\
cc65 C & \cmd{make lib}; include \cmd{ultimate.h} and link
  \cmd{bindings/cc65/build/ultimate.lib} \\
ca65 assembly & link the same library and include
  \cmd{bindings/asm/ultimate.inc} \\
other toolchains & \cmd{make blob}; load the standalone binary described in
  Chapter~5 \\
\end{reftable}

\evidence{\cmd{Makefile}; \cmd{src/basic/Makefile};
\cmd{bindings/cc65/Makefile}; \cmd{bindings/blob/Makefile}}

\Needspace{18\baselineskip}
\subsection*{BASIC check}

Load and run the wedge, list the current directory, then print the command
result:

\begin{screen}
LOAD"UCI",8\\
RUN\\
UDIR\\
PRINT UERR\\
0
\end{screen}

This demonstrates wedge installation, a DOS command and the shared result
model. A printed \bas{0} is \cmd{ULTIMATE\_OK}; \bas{1} is
\cmd{ULTIMATE\_ERR\_NO\_DEVICE}.

\evidence{\cmd{src/basic/loader.s}; \cmd{src/basic/dispatch.s};
\cmd{include/uci\_protocol.h}}

\subsection*{C check}

This complete program demonstrates initialization and error reporting:

\begin{listingbox}
\#include <stdio.h>\\
\#include <ultimate.h>\\
\\
int main(void)\\
\{\\
\hspace*{1em}uint8\_t err = ultimate\_init();\\
\\
\hspace*{1em}if (err != ULTIMATE\_OK)\\
\hspace*{2em}printf("ultimate: \%s\textbackslash n", ultimate\_strerror(err));\\
\hspace*{1em}return err != ULTIMATE\_OK;\\
\}
\end{listingbox}

\begin{outputbox}
No text on success. On failure, for example:\\
ULTIMATE: NO ULTIMATE FOUND
\end{outputbox}

Build it from the repository root:

\begin{listingbox}
cl65 -t c64 -I include -o hello.prg hello.c \textbackslash\\
\hspace*{2em}bindings/cc65/build/ultimate.lib
\end{listingbox}

\begin{outputbox}
hello.prg is created with no compiler errors.
\end{outputbox}

\evidence{\cmd{include/ultimate.h}; \cmd{examples/cc65/Makefile}}

\Needspace{14\baselineskip}
\subsection*{Assembly check}

This demonstrates the assembly return convention: the result code is in
\cmd{A}, ready to compare and branch on.

\begin{listingbox}
\hspace*{2em}jsr ultimate\_init\\
\hspace*{2em}cmp \#ULTIMATE\_OK\\
\hspace*{2em}bne failed
\end{listingbox}

\begin{outputbox}
A = ULTIMATE_OK; BNE is not taken.
\end{outputbox}

The \cmd{cmp} instruction leaves the processor flags ready for \cmd{bne}.

\evidence{\cmd{docs/asm-abi.md}; \cmd{examples/asm/identify.s}}

\section{Recover From Common Failures}

\subsection*{No device}

\cmd{ULTIMATE\_ERR\_NO\_DEVICE} means the signature is absent. Check the
Command Interface setting before assuming the hardware is missing.

\Needspace{12\baselineskip}
\subsection*{An abandoned exchange}

A program can stop while the firmware is holding command or reply state.
\cmd{uci\_abort()} releases it; \cmd{ultimate\_init()} performs that recovery
during startup. From BASIC, a bare \bas{uci} calls the same abort path:

\begin{screen}
UCI:PRINT UERR\\
0\\
READY.
\end{screen}

This example demonstrates recovery only; it does not send a command.

\subsection*{A missing feature}

Use \cmd{ultimate\_detect()} for command-interface targets,
\cmd{ultimate\_reu\_available()} for the REU, and
\cmd{ultimate\_turbo\_available()} for turbo. Do not infer capabilities from a
model name.

\subsection*{A short buffer}

\cmd{ULTIMATE\_ERR\_TRUNCATED} means an exchange completed but not all reply
bytes fitted. Output lengths report the bytes retained, not the bytes discarded.

\evidence{\cmd{include/uci.h}; \cmd{include/ultimate.h};
\cmd{src/basic/dispatch.s}; \cmd{tests/emulator/sdk.suite}}

% Chapter 2 - BASIC.

\guidechapter{PROGRAMMING IN BASIC}{%
  \item Installing the wedge
  \item The twenty-three keywords
  \item Files and directories
  \item The RAM expansion
  \item Turbo
  \item Talking to the Ultimate directly
  \item What BASIC cannot do
}

\section{Installing The Wedge}

The wedge comes two ways, and they are the same program.

\textbf{As a program.} \bas{load"uci",8} then \bas{run}. It installs itself
into the 4K block at \cmd{\$C000}, which BASIC does not use, and leaves your
38911 bytes free exactly as they were.

\textbf{As a cartridge.} Mount \cmd{uci.crt} from the Ultimate's menu and the
wedge is there the moment the machine starts, with nothing to load. The
cartridge costs BASIC its 8K: you get 30719 bytes free instead of 38911.

Either way the machine greets you and the keywords are live. They survive
\bas{new}, they survive \bas{run/stop}~+~\key{restore}, and they survive a
program crashing. They do not survive a reset.

\begin{note}
The wedge takes over four of BASIC's own vectors --- the ones for tokenising,
listing, statement dispatch and expression evaluation. Every other wedge and
fast-loader that does the same thing will fight with it. If you use one, load
the SDK's wedge last.
\end{note}

\section{The Twenty-Three Keywords}

They fall into four groups, and you can ignore three of them until you need
them.

\subsection*{Asking how things went}

Nothing in this wedge stops your program with an error. A command that fails
sets a result code and returns, because a demo dropping to \bas{ready.} in the
middle of a part is worse than one that quietly did nothing. So every command
is followed, when it matters, by a look at \bas{uerr}.

\begin{reftable}{lp{0.62\linewidth}}
\bas{uerr} & the result of the last command. \bas{0} is success \\
\bas{udev} & the raw code the Ultimate reported, for diagnosis \\
\bas{ulen} & how many bytes came back \\
\bas{ust\$} & the status text, exactly as the Ultimate sent it \\
\bas{udat\$} & the reply as a string, up to 255 bytes \\
\bas{ubyte(n)} & byte \bas{n} of the reply, or \bas{0} past the end \\
\end{reftable}

\subsection*{Files, and the RAM expansion}

\begin{reftable}{lp{0.62\linewidth}}
\bas{udir} & print the current directory \\
\bas{uload} \textit{name\$} [,\textit{addr}] & load a PRG; no address means the
  file's own \\
\bas{ubload} \textit{name\$}, \textit{addr}, \textit{len} & load raw bytes, and
  not one past \textit{len} \\
\bas{usave} \textit{name\$}, \textit{addr}, \textit{len} & memory to a file \\
\bas{ustash} \textit{addr}, \textit{reu}, \textit{len} & C64 memory into the
  expansion \\
\bas{ufetch} \textit{addr}, \textit{reu}, \textit{len} & and back out of it \\
\bas{ureu} & the expansion's size in 64K banks, \bas{0} for none \\
\end{reftable}

\subsection*{Speed}

\begin{reftable}{lp{0.62\linewidth}}
\bas{uturbo} \textit{n} & set the processor speed index \\
\bas{uturbo} & as a function, the current index --- or \bas{255} when the
  machine's owner has turned turbo off \\
\end{reftable}

\subsection*{Everything else}

\begin{reftable}{lp{0.62\linewidth}}
\bas{uci} \textit{t}, \textit{c} [, \textit{args}] & send any command to any
  target \\
\bas{uci} & on its own, rescue a wedged interface \\
\bas{uw(}\textit{x}\bas{)} & force a 16-bit argument \\
\bas{ul(}\textit{x}\bas{)} & force a 32-bit argument \\
\bas{udos1} \bas{udos2} \bas{unet} \bas{uctrl} \bas{uiec} \bas{uhttp} &
  the six target numbers, by name \\
\end{reftable}

\section{Files And Directories}

\bas{udir} prints the directory and is the quickest way to see that everything
is working:

\begin{screen}
UDIR
\end{screen}

Loading is the interesting one, because the Ultimate has a fast path and the
SDK takes it when it can. This loads a file to the address stored in the file
itself, exactly as \bas{load"x",8,1} would:

\begin{screen}
10 ULOAD "GAME"\\
20 IF UERR THEN PRINT "COULD NOT LOAD": END
\end{screen}

Give it an address and it loads there instead:

\begin{screen}
10 ULOAD "SPRITES",832\\
20 IF UERR THEN PRINT "NO SPRITES": END
\end{screen}

\bas{ubload} is the one to use for data rather than programs. It interprets
nothing, strips no header, and stops at the length you give it:

\begin{screen}
10 UBLOAD "MUSIC.BIN",49152,4096
\end{screen}

\begin{note}
Filenames go to the Ultimate exactly as you type them. Uppercase PETSCII and
uppercase ASCII are the same bytes, so a name typed at a C64 arrives as the
Ultimate expects it, and the filesystem does not care about case anyway.
\end{note}

\section{The RAM Expansion}

If there is an expansion fitted, you have up to sixteen megabytes to put things
in --- and moving bytes in and out of it costs almost nothing, because the REU
does it with the processor halted rather than a byte at a time.

Always ask first. \bas{ureu} returns the size in 64K banks and zero when there
is nothing there, so one test covers both questions:

\begin{screen}
10 B = UREU\\
20 IF B = 0 THEN PRINT "NO EXPANSION": END\\
30 PRINT "EXPANSION IS"; B * 64; "K"
\end{screen}

Now put something in it and get it back. This copies the top of the screen into
the expansion, clears the screen, and puts it back:

\begin{screen}
10 IF UREU = 0 THEN PRINT "NEED AN REU": END\\
20 USTASH 1024,0,1000\\
30 PRINT CHR\$(147)\\
40 FOR I = 1 TO 2000: NEXT\\
50 UFETCH 1024,0,1000
\end{screen}

Line~20 sends 1000 bytes from screen memory at 1024 to expansion address~0.
Line~50 brings them back. Both are instant.

\begin{note}
The expansion address is a full 24-bit number, so it can reach all sixteen
megabytes. BASIC handles that fine --- \bas{ustash 1024, 1048576, 256} is a
perfectly good line.
\end{note}

\begin{warning}
Writing past the end of the expansion does not fail --- it wraps round to the
beginning and quietly overwrites what is there. That is exactly how \bas{ureu}
measures the size, and it is why asking first is worth the one line.
\end{warning}

\section{Turbo}

On an Ultimate~64 the processor can run faster than 1MHz, and \bas{uturbo}
sets it. The catch is that this is the machine owner's decision, not your
program's: unless \bas{turbo control} is set in the Ultimate's menu, the
registers are not there at all.

\begin{screen}
10 UTURBO 3\\
20 IF UERR THEN PRINT "NO TURBO HERE"\\
30 PRINT "SPEED INDEX IS"; UTURBO
\end{screen}

Index \bas{0} is 1MHz and index \bas{3} is 4MHz on every machine that has
turbo. Above~3 the numbers mean different speeds on different machines, so a
program that ships to other people should stay at or below~3 or ask.

The function form answers \bas{255} when turbo is unavailable. That is not a
speed --- the index is only four bits --- so you can test for it.

\section{Talking To The Ultimate Directly}

The wedge wraps the operations most programs want. Everything else --- and
there are around a hundred commands --- is reachable with \bas{uci}, which
sends any command to any target.

\begin{screen}
10 UCI UCTRL, 1\\
20 PRINT UDAT\$
\end{screen}

That sends command~1 to the control target, which identifies itself, and prints
what came back.

Arguments follow the command, and the wedge sends each number as a single byte
because that is what nearly every command wants. When a command wants something
wider, say so with \bas{uw()} for sixteen bits or \bas{ul()} for thirty-two:

\begin{screen}
10 UCI UDOS1, 16, UW(8192)
\end{screen}

\begin{note}
The SDK knows the shape of many commands and will widen arguments for you
without being told. \bas{uw()} and \bas{ul()} are the escape hatch for a
command newer than the SDK's table, and they always win over it.
\end{note}

\section{What BASIC Cannot Do}

Two limits worth knowing before they surprise you.

\subsection*{IF X THEN UCI does not work}

BASIC's \bas{if} jumps straight into the ROM's own statement dispatcher, which
rejects any token above \bas{new}. Every extension that adds statement keywords
to BASIC~V2 has this hole. So this does not work:

\begin{screen}
10 IF X = 1 THEN ULOAD "GAME"
\end{screen}

and this does:

\begin{screen}
10 IF X = 1 THEN 100\\
\hspace*{1em}\ldots\\
100 ULOAD "GAME"
\end{screen}

Functions are unaffected, because expression evaluation does come back through
the wedge. \bas{if uerr then \ldots} is fine, and so is \bas{if ureu = 0 then
\ldots}.

\subsection*{Strings are 255 bytes}

\bas{udat\$} gives you the reply as a string, clipped at 255 bytes because that
is the longest string BASIC has. When a reply is longer than that, use
\bas{ubyte(n)} to walk it, or \bas{ubload} to put it somewhere you can
\bas{peek}.

% Chapter 3 - C.

\guidechapter{PROGRAMMING IN C}{%
  \item Building against the SDK
  \item The shape of every call
  \item Buffers belong to you
  \item Result codes
  \item Capability detection
  \item A worked example
}

\section{Building Against The SDK}

Build the library once:

\begin{listingbox}
make lib
\end{listingbox}

then compile against \cmd{include/} and link \cmd{ultimate.lib}:

\begin{listingbox}
cl65 -t c64 -Osir -I include -o prog.prg prog.c \textbackslash\\
\hspace*{2em}bindings/cc65/build/ultimate.lib
\end{listingbox}

One header covers the whole library:

\begin{listingbox}
\#include <ultimate.h>
\end{listingbox}

If you need the transport itself --- to send a command the SDK does not wrap ---
add \cmd{<uci.h>} as well. Most programs never do.

\section{The Shape Of Every Call}

Almost every function in the SDK looks like this:

\begin{listingbox}
uint8\_t ultimate\_something(\ldots);
\end{listingbox}

It returns a result code, and \cmd{ULTIMATE\_OK} is zero. Anything it produces
comes back through a pointer you passed in. There are no global error variables
to check, nothing is returned in a static buffer, and no call allocates memory.

\begin{listingbox}
uint16\_t got;\\
uint8\_t  err;\\
\\
err = ultimate\_read(buf, sizeof buf, \&got);\\
if (err != ULTIMATE\_OK)\\
\hspace*{1em}return err;
\end{listingbox}

\section{Buffers Belong To You}

The SDK has no heap and no hidden buffers. Every byte it produces lands in
memory you own, and it never keeps a pointer after the call returns. That is
what makes it usable from an interrupt-driven demo, and it is why every call
that produces data takes both a buffer and its size.

\begin{warning}
The size you pass is the size of the buffer, not the number of bytes you want.
If a reply is larger, you get \cmd{ULTIMATE\_ERR\_TRUNCATED}, the count tells
you how much was kept, and nothing is written past the end.
\end{warning}

\section{Result Codes}

There are eleven, they never change, and Appendix~B lists them. Three are worth
knowing straight away:

\begin{reftable}{lp{0.60\linewidth}}
\cmd{ULTIMATE\_OK} & it worked \\
\cmd{ULTIMATE\_END} & there is no more --- the end of a directory, or a socket
  the far end has closed. Not a failure \\
\cmd{ULTIMATE\_ERR\_TRUNCATED} & it worked, and your buffer was too small \\
\end{reftable}

Turn any of them into text with:

\begin{listingbox}
printf("\%s\textbackslash n", ultimate\_strerror(err));
\end{listingbox}

\section{Capability Detection}

Different machines have different targets, and firmware versions differ too.
Ask once, at start-up, rather than assuming:

\begin{listingbox}
ultimate\_capabilities caps;\\
\\
ultimate\_detect(\&caps);\\
if (ultimate\_has\_http(\&caps)) \{ /* fetch things */ \}\\
if (ultimate\_has\_network(\&caps)) \{ /* sockets */ \}
\end{listingbox}

The hardware that is not reached through the command interface answers
separately, because it is not a target: \cmd{ultimate\_turbo\_available()} and
\cmd{ultimate\_reu\_size()} tell you about the processor speed registers and
the RAM expansion.

\section{A Worked Example}

This prints the directory, one entry at a time. It is the pattern for every
walk in the SDK: start it, take entries until \cmd{ULTIMATE\_END}, and check
the result of the call that ended it.

\begin{listingbox}
\#include <stdio.h>\\
\#include <ultimate.h>\\
\\
int main(void)\\
\{\\
\hspace*{1em}char    name[64];\\
\hspace*{1em}uint8\_t attrib;\\
\hspace*{1em}uint8\_t err;\\
\\
\hspace*{1em}if (ultimate\_init() != ULTIMATE\_OK) \{\\
\hspace*{2em}printf("no ultimate\textbackslash n");\\
\hspace*{2em}return 1;\\
\hspace*{1em}\}\\
\\
\hspace*{1em}err = ultimate\_opendir();\\
\hspace*{1em}while (err == ULTIMATE\_OK) \{\\
\hspace*{2em}err = ultimate\_readdir(name, sizeof name, \&attrib);\\
\hspace*{2em}if (err != ULTIMATE\_OK)\\
\hspace*{3em}break;\\
\hspace*{2em}printf("\%s\%s\textbackslash n", name,\\
\hspace*{3em}(attrib \& DOS\_ATTR\_DIR) ? "/" : "");\\
\hspace*{1em}\}\\
\\
\hspace*{1em}if (err != ULTIMATE\_END)\\
\hspace*{2em}printf("stopped: \%s\textbackslash n", ultimate\_strerror(err));\\
\hspace*{1em}return 0;\\
\}
\end{listingbox}

\begin{note}
A directory walk is one live exchange with the Ultimate. Do not issue any other
command in the middle of one --- the block boundaries between entries are the
only thing separating one name from the next. If you need to stop early, call
\cmd{uci\_abort()}.
\end{note}

% Chapter 4 - Assembly.

\guidechapter{PROGRAMMING IN ASSEMBLY}{%
  \item The register and shared-variable conventions
  \item Relocating the SDK's RAM and zero page
  \item Every service entry point by example
  \item Raw request blocks with optional spans
}

\section{The Calling Convention}

Include \cmd{ultimate.inc} and call the unprefixed entry points. Unless an
entry below says otherwise, a service returns an \cmd{ULTIMATE\_*} result in
\cmd{A} and may change \cmd{A}, \cmd{X}, \cmd{Y} and the flags. A byte argument
uses \cmd{A}; a pointer argument uses low byte in \cmd{A} and high byte in
\cmd{X}; multi-byte service arguments live in the exported \cmd{ult\_*}
variables.

The examples use these two ordinary ca65 macros to keep pointer and word setup
visible without repeating four instructions each time:

\begin{listingbox}
\hspace*{2em}.include "ultimate.inc"\
\
.macro setptr location, value\
\hspace*{2em}lda \#<value\
\hspace*{2em}sta location\
\hspace*{2em}lda \#>value\
\hspace*{2em}sta location + 1\
.endmacro\
\
.macro setword location, value\
\hspace*{2em}lda \#<value\
\hspace*{2em}sta location\
\hspace*{2em}lda \#>value\
\hspace*{2em}sta location + 1\
.endmacro
\end{listingbox}

\begin{outputbox}
setptr ult_buf, text stores the address of text in ult_buf.
\end{outputbox}

\cmd{setptr} and \cmd{setword} emit the same bytes; their different names make
the purpose clear at each call site. Result branches always compare \cmd{A}
immediately, before another instruction replaces it.

Assembly output boxes show the returned register or memory values from a
representative successful run. They omit screen-printing routines so the SDK
operation itself stays visible.

With the C64 character map, lowercase source letters assemble to the uppercase
byte values shared by PETSCII and ASCII. Paths, filenames and HTTP text
therefore appear in lowercase source below but reach the firmware in uppercase.

\evidence{\cmd{bindings/asm/ultimate.inc}; \cmd{docs/asm-abi.md}}

\section{Place The SDK's Working Memory}

The SDK's zero-page block starts at \cmd{UCI\_ZP}, which defaults to
\cmd{\$FB}. Its size is \cmd{UCI\_ZP\_SIZE}. The normal variable block occupies
\cmd{UCI\_VARS\_SIZE} bytes in the BSS segment, which is where the linker puts
variables that have no starting value. Define both bases when assembling every
SDK source to choose fixed locations instead:

\begin{listingbox}
ca65 -D UCI\_VARS=\$CF00 -D UCI\_ZP=\$A3 uci\_core.s
\end{listingbox}

\begin{outputbox}
UCI variables begin at \$CF00; zero-page workspace begins at \$A3.
\end{outputbox}

With \cmd{UCI\_VARS} defined, the core emits no BSS; its variables become
equates from that address. The SDK installs no interrupt handler and is not
re-entrant, so an interrupt routine must not call it while foreground code is
inside a service.

\evidence{\cmd{bindings/asm/uci\_protocol.inc};
\cmd{src/uci/uci\_zp.inc}; \cmd{docs/asm-abi.md}}

\Needspace{30\baselineskip}
\section{Initialize, Detect And Describe The Machine}

\begin{listingbox}
\hspace*{2em}jsr ultimate\_init\
\hspace*{2em}cmp \#ULTIMATE\_OK\
\hspace*{2em}bne failed\
\
\hspace*{2em}jsr ultimate\_available\
\hspace*{2em}beq unavailable\
\
\hspace*{2em}lda \#<capabilities\
\hspace*{2em}ldx \#>capabilities\
\hspace*{2em}jsr ultimate\_detect\
\hspace*{2em}cmp \#ULTIMATE\_OK\
\hspace*{2em}bne failed\
\
capabilities: .res 4
\end{listingbox}

\begin{outputbox}
A = ULTIMATE_OK\\
capabilities records the detected targets.
\end{outputbox}

\cmd{ultimate\_init} performs startup. \cmd{ultimate\_available} returns a
boolean rather than a result code. \cmd{ultimate\_detect} takes the address of
the four-byte capability block directly in \cmd{A}/\cmd{X}.

\Needspace{20\baselineskip}
\subsection*{Identify a target with and without the optional length output}

\begin{listingbox}
\hspace*{2em}setptr ult\_buf, text\
\hspace*{2em}setword ult\_buflen, 64\
\hspace*{2em}setptr ult\_outlen, text\_length\
\hspace*{2em}lda \#UCI\_TARGET\_DOS1\
\hspace*{2em}jsr ultimate\_identify\
\
\hspace*{2em}setword ult\_outlen, 0\
\hspace*{2em}lda \#UCI\_TARGET\_DOS1\
\hspace*{2em}jsr ultimate\_identify\
\
text:        .res 64\
text\_length: .res 2
\end{listingbox}

\begin{outputbox}
text = "ULTIMATE-II DOS V1.2"\\
text_length = 20 with the pointer; no length is stored with zero.
\end{outputbox}

Both calls terminate the text in \cmd{ult\_buf}. The first stores its length
through \cmd{ult\_outlen}; zero disables that optional output.

\Needspace{25\baselineskip}
\subsection*{Read the model with and without the optional length output}

\begin{listingbox}
\hspace*{2em}setptr ult\_buf, text\
\hspace*{2em}setword ult\_buflen, 64\
\hspace*{2em}setptr ult\_outlen, text\_length\
\hspace*{2em}jsr ultimate\_get\_model\
\
\hspace*{2em}setword ult\_outlen, 0\
\hspace*{2em}jsr ultimate\_get\_model
\end{listingbox}

\begin{outputbox}
text = "Ultimate 64 Elite"\\
text_length is optional in the same way as IDENTIFY.
\end{outputbox}

\cmd{ultimate\_get\_model} uses the same buffer contract but needs no target
argument. Its underlying \cmd{CTRL\_CMD\_GET\_HWINFO} command is marked
deprecated by the firmware documentation; this existing wrapper remains for
compatibility.

\Needspace{32\baselineskip}
\subsection*{Read SID addresses through deprecated HWINFO}

\begin{listingbox}
\hspace*{2em}lda \#<sid\_info\
\hspace*{2em}ldx \#>sid\_info\
\hspace*{2em}jsr ultimate\_legacy\_get\_sid\_info\
\hspace*{2em}cmp \#ULTIMATE\_OK\
\hspace*{2em}bne failed\
\
\hspace*{2em}lda sid\_info + ULTIMATE\_SID\_INFO\_COUNT\
\hspace*{2em}; iterate this many ULTIMATE_SID_RECORD_SIZE-byte records\
\hspace*{2em}ldx \#ULTIMATE\_SID\_INFO\_RECORDS\
\hspace*{2em}lda sid\_info + ULTIMATE\_SID\_PRIMARY\_ADDRESS,x\
\hspace*{2em}sta first\_sid\_address\
\hspace*{2em}lda sid\_info + ULTIMATE\_SID\_PRIMARY\_ADDRESS + 1,x\
\hspace*{2em}sta first\_sid\_address + 1\
\
sid\_info:          .res ULTIMATE\_SID\_INFO\_SIZE\
first\_sid\_address: .res 2
\end{listingbox}

\begin{outputbox}
sid_info count = 4\\
primary addresses = \$D400, \$D400, \$D400, \$D400
\end{outputbox}

Records begin at \cmd{ULTIMATE\_SID\_INFO\_RECORDS}. Within each record,
\cmd{ULTIMATE\_SID\_PRIMARY\_ADDRESS} and
\cmd{ULTIMATE\_SID\_SECONDARY\_ADDRESS} name two little-endian words, and
\cmd{ULTIMATE\_SID\_TYPE} names the raw type byte. The type is not a 6581 or
8580 identifier. Firmware still implements this query, but marks HWINFO
deprecated and publishes no replacement C64-side UCI command. The legacy name
is deliberate; handle \cmd{ULTIMATE\_ERR\_NOT\_SUPPORTED} and read no record
unless the call returned \cmd{ULTIMATE\_OK}.

\evidence{\cmd{bindings/asm/ultimate.inc}; \cmd{src/uci/ultimate.s};
\cmd{software/io/command\_interface/control\_target.cc}}

\Needspace{15\baselineskip}
\subsection*{Turn a result into text}

\begin{listingbox}
\hspace*{2em}lda error\_code\
\hspace*{2em}jsr ultimate\_strerror\
\hspace*{2em}sta message\_ptr\
\hspace*{2em}stx message\_ptr + 1\
\
error\_code: .res 1\
message\_ptr: .res 2
\end{listingbox}

\begin{outputbox}
message_ptr points to the PETSCII text for error_code.
\end{outputbox}

\cmd{ultimate\_strerror} returns the PETSCII message pointer in
\cmd{A}/\cmd{X}; it does not return a result code.

\evidence{\cmd{docs/asm-abi.md}; \cmd{src/uci/ultimate.s};
\cmd{src/uci/ultimate\_strerror.s}}

\Needspace{32\baselineskip}
\section{Read And Change The Palette}

\begin{listingbox}
\hspace*{2em}lda \#<palette\
\hspace*{2em}ldx \#>palette\
\hspace*{2em}jsr ultimate\_palette\_get\
\
\hspace*{2em}lda \#<palette\
\hspace*{2em}ldx \#>palette\
\hspace*{2em}jsr ultimate\_palette\_set\
\
\hspace*{2em}lda \#6\
\hspace*{2em}sta ult\_color\
\hspace*{2em}lda \#255\
\hspace*{2em}sta ult\_color + 1\
\hspace*{2em}lda \#0\
\hspace*{2em}sta ult\_color + 2\
\hspace*{2em}sta ult\_color + 3\
\hspace*{2em}jsr ultimate\_palette\_set\_color\
\
\hspace*{2em}jsr ultimate\_palette\_reset\
\
palette: .res UCI\_PALETTE\_BYTES
\end{listingbox}

\begin{outputbox}
GET fills 48 bytes; SET COLOUR 6 TO RED and RESET return ULTIMATE_OK.
\end{outputbox}

Get and set take a direct pointer in \cmd{A}/\cmd{X}. The single-colour entry
reads index, red, green and blue from the four consecutive \cmd{ult\_color}
bytes. Reset restores the built-in live palette.
Firmware without these commands returns
\cmd{ULTIMATE\_ERR\_NOT\_SUPPORTED}; branch on \cmd{A} before using the
48-byte result from \cmd{ultimate\_palette\_get}.

\evidence{\cmd{docs/asm-abi.md}; \cmd{src/uci/palette.s}}

\section{Use Turbo And Badline Control}

\begin{listingbox}
\hspace*{2em}jsr ultimate\_turbo\_available\
\hspace*{2em}beq no\_turbo\
\
\hspace*{2em}jsr ultimate\_turbo\_get\
\hspace*{2em}sta old\_speed\
\
\hspace*{2em}lda \#U64\_SPEED\_4MHZ\
\hspace*{2em}jsr ultimate\_turbo\_set\
\
\hspace*{2em}lda \#0\
\hspace*{2em}jsr ultimate\_turbo\_badlines\
\hspace*{2em}lda \#1\
\hspace*{2em}jsr ultimate\_turbo\_badlines\
\
\hspace*{2em}lda old\_speed\
\hspace*{2em}jsr ultimate\_turbo\_set\
\
old\_speed: .res 1
\end{listingbox}

\begin{outputbox}
old_speed = previous speed index\\
SET 4 MHZ, THEN RESTORE SPEED AND BADLINES: ULTIMATE_OK
\end{outputbox}

The available and get calls return values directly. Set returns a result. The
two badline calls show both forms explicitly: zero disables them; non-zero
restores them. The final set restores the speed read at entry.

\evidence{\cmd{docs/asm-abi.md}; \cmd{src/uci/turbo.s}}

\section{Change And Read The Current Directory}

\begin{listingbox}
\hspace*{2em}setptr ult\_buf, path\
\hspace*{2em}jsr ultimate\_chdir\
\
\hspace*{2em}setptr ult\_buf, path\_buffer\
\hspace*{2em}setword ult\_buflen, 128\
\hspace*{2em}setptr ult\_outlen, path\_length\
\hspace*{2em}jsr ultimate\_getpath\
\
\hspace*{2em}setword ult\_outlen, 0\
\hspace*{2em}jsr ultimate\_getpath\
\
path:        .byte "/usb1/games", 0\
path\_buffer: .res 128\
path\_length: .res 2
\end{listingbox}

\begin{outputbox}
path_buffer = "/USB1/GAMES"\\
path_length = 11 when the optional pointer is enabled.
\end{outputbox}

\cmd{ultimate\_chdir} consumes the terminated path through \cmd{ult\_buf}.
The two getpath calls show the optional length pointer enabled and disabled.

\subsection*{Read every directory entry}

\begin{listingbox}
\hspace*{2em}jsr ultimate\_opendir\
\hspace*{2em}cmp \#ULTIMATE\_OK\
\hspace*{2em}bne failed\
\
next\_entry:\
\hspace*{2em}setptr ult\_buf, name\
\hspace*{2em}setword ult\_buflen, 64\
\hspace*{2em}jsr ultimate\_readdir\
\hspace*{2em}cmp \#ULTIMATE\_END\
\hspace*{2em}beq directory\_done\
\hspace*{2em}cmp \#ULTIMATE\_OK\
\hspace*{2em}bne failed\
\hspace*{2em}lda ult\_attrib\
\hspace*{2em}; use name and its DOS\_ATTR\_* bits here\
\hspace*{2em}jmp next\_entry\
\
name: .res 64
\end{listingbox}

\begin{outputbox}
name receives one entry per call; A becomes ULTIMATE_END after the last.
\end{outputbox}

Each readdir releases one reply block and writes attributes to
\cmd{ult\_attrib}. Issue no other command inside the walk. Call
\cmd{uci\_abort} if leaving before \cmd{ULTIMATE\_END}.

\subsection*{Make a directory, and return to the configured one}

\begin{listingbox}
\hspace*{2em}setptr ult\_buf, dir\_name\
\hspace*{2em}jsr ultimate\_mkdir\
\
\hspace*{2em}jsr ultimate\_home\
\
dir\_name: .byte "saves", 0
\end{listingbox}

\begin{outputbox}
SAVES is created; A = ULTIMATE_OK from each call.
\end{outputbox}

\cmd{ultimate\_mkdir} takes the terminated name through \cmd{ult\_buf} and
measures it itself. \cmd{ultimate\_home} takes nothing at all.

\evidence{\cmd{docs/asm-abi.md}; \cmd{src/uci/dos.s}}

\section{Open, Seek, Read, Write And Close Files}

\begin{listingbox}
\hspace*{2em}setptr ult\_buf, read\_name\
\hspace*{2em}lda \#DOS\_FA\_READ\
\hspace*{2em}jsr ultimate\_open\
\hspace*{2em}cmp \#ULTIMATE\_OK\
\hspace*{2em}bne failed\
\
\hspace*{2em}lda \#0\
\hspace*{2em}sta ult\_num\
\hspace*{2em}sta ult\_num + 1\
\hspace*{2em}sta ult\_num + 2\
\hspace*{2em}sta ult\_num + 3\
\hspace*{2em}jsr ultimate\_seek\
\
\hspace*{2em}setptr ult\_buf, file\_buffer\
\hspace*{2em}setword ult\_buflen, 128\
\hspace*{2em}setptr ult\_outlen, bytes\_read\
\hspace*{2em}jsr ultimate\_read\
\hspace*{2em}jsr ultimate\_close\
\
read\_name:   .byte "readme.txt", 0\
file\_buffer: .res 128\
bytes\_read:  .res 2
\end{listingbox}

\begin{outputbox}
file_buffer receives up to 128 bytes; bytes_read receives the actual count.
\end{outputbox}

Open establishes the firmware's current file. Seek selects its absolute
position. Read stores bytes and, when \cmd{ult\_outlen} is non-zero, their
count. Close releases that same firmware-side file.

\Needspace{13\baselineskip}
The count output is optional:

\begin{listingbox}
\hspace*{2em}; With a readable file already open:\
\hspace*{2em}setptr ult\_buf, file\_buffer\
\hspace*{2em}setword ult\_buflen, 128\
\hspace*{2em}setword ult\_outlen, 0\
\hspace*{2em}jsr ultimate\_read
\end{listingbox}

\begin{outputbox}
file_buffer receives bytes; no byte count is stored.
\end{outputbox}

\subsection*{Create, write and delete a file}

\begin{listingbox}
\hspace*{2em}setptr ult\_buf, write\_name\
\hspace*{2em}lda \#DOS\_FA\_CREATE\_ALWAYS | DOS\_FA\_WRITE\
\hspace*{2em}jsr ultimate\_open\
\hspace*{2em}cmp \#ULTIMATE\_OK\
\hspace*{2em}bne failed\
\
\hspace*{2em}setptr ult\_buf, message\
\hspace*{2em}setword ult\_buflen, message\_end - message\
\hspace*{2em}jsr ultimate\_write\
\hspace*{2em}jsr ultimate\_close\
\
\hspace*{2em}setptr ult\_buf, write\_name\
\hspace*{2em}jsr ultimate\_delete\
\
write\_name: .byte "result.txt", 0\
message:    .byte "ok", 13\
message\_end:
\end{listingbox}

\begin{outputbox}
RESULT.TXT: 3 bytes written, closed and deleted; A = ULTIMATE_OK.
\end{outputbox}

The open mask creates or replaces the file. Write consumes
\cmd{ult\_buflen} bytes from \cmd{ult\_buf}; delete consumes the terminated
name after the file is closed.

\evidence{\cmd{docs/asm-abi.md}; \cmd{src/uci/dos.s}}

\Needspace{28\baselineskip}
\section{Ask What A File Is}

\begin{listingbox}
\hspace*{2em}setptr ult\_buf, stat\_name\
\hspace*{2em}setptr ult\_arg2, file\_info\
\hspace*{2em}jsr ultimate\_stat\
\hspace*{2em}cmp \#ULTIMATE\_OK\
\hspace*{2em}bne failed\
\
\hspace*{2em}lda file\_info + DOS\_INFO\_SIZE\
\hspace*{2em}sta file\_size\
\hspace*{2em}lda file\_info + DOS\_INFO\_SIZE + 1\
\hspace*{2em}sta file\_size + 1\
\hspace*{2em}lda file\_info + DOS\_INFO\_ATTRIB\
\hspace*{2em}and \#DOS\_ATTR\_DIR\
\hspace*{2em}bne is\_a\_directory\
\
stat\_name: .byte "readme.txt", 0\
file\_info: .res ULTIMATE\_FILEINFO\_SIZE\
file\_size: .res 2
\end{listingbox}

\begin{outputbox}
file_info holds the reply; the name inside it is NUL-terminated.
\end{outputbox}

\cmd{ult\_buf} is the name and \cmd{ult\_arg2} is the block the reply is
received into. That block must be \cmd{ULTIMATE\_FILEINFO\_SIZE} bytes.
\cmd{ult\_buflen} is not an input here: \cmd{ultimate\_stat} measures the name
itself and overwrites \cmd{ult\_buflen} doing it.

\cmd{ultimate\_fstat} asks the same question about the file
\cmd{ultimate\_open} left open, so it takes only the block:

\begin{listingbox}
\hspace*{2em}setptr ult\_arg2, file\_info\
\hspace*{2em}jsr ultimate\_fstat
\end{listingbox}

\begin{outputbox}
file_info describes the open file; ult_buf is not read.
\end{outputbox}

The \cmd{DOS\_INFO\_*} equates in \cmd{uci\_protocol.inc} name the fields
inside the block:

\begin{reftable}{lp{0.60\linewidth}}
\cmd{DOS\_INFO\_SIZE} & \cmd{+\$00}, size in bytes, 32-bit, low byte first \\
\cmd{DOS\_INFO\_DATE} & \cmd{+\$04}, year less 1980 in bits 15--9, month 8--5, day 4--0 \\
\cmd{DOS\_INFO\_TIME} & \cmd{+\$06}, hour 15--11, minute 10--5, two-second units 4--0 \\
\cmd{DOS\_INFO\_EXTENSION} & \cmd{+\$08}, three bytes, space padded \\
\cmd{DOS\_INFO\_ATTRIB} & \cmd{+\$0B}, the \cmd{DOS\_ATTR\_*} bits \\
\cmd{DOS\_INFO\_NAME} & \cmd{+\$0C}, the name, terminated by the SDK \\
\end{reftable}

Date and time are the packed forms used by FAT, the File Allocation Table
filesystem the Ultimate's storage uses, and they arrive unconverted.

\evidence{\cmd{bindings/asm/ultimate.inc}; \cmd{src/uci/dosinfo.s};
\cmd{bindings/asm/uci\_protocol.inc}}

\section{Rename And Copy A File}

\begin{listingbox}
\hspace*{2em}setptr ult\_buf, old\_name\
\hspace*{2em}setptr ult\_arg2, new\_name\
\hspace*{2em}jsr ultimate\_rename\
\
\hspace*{2em}setptr ult\_buf, new\_name\
\hspace*{2em}setptr ult\_arg2, copy\_name\
\hspace*{2em}jsr ultimate\_copy\
\
old\_name:  .byte "old.txt", 0\
new\_name:  .byte "new.txt", 0\
copy\_name: .byte "backup.txt", 0
\end{listingbox}

\begin{outputbox}
OLD.TXT becomes NEW.TXT; BACKUP.TXT is a copy of it.
\end{outputbox}

Both take the first name in \cmd{ult\_buf} and the second in \cmd{ult\_arg2}.
The request block carries two spans and these commands put three runs of bytes
on the wire, so the first name's own terminator separates it from the second.
The Ultimate performs the copy, so no bytes pass through the C64.

\evidence{\cmd{bindings/asm/ultimate.inc}; \cmd{src/uci/dosinfo.s}}

\section{Load And Save C64 Memory}

\begin{listingbox}
\hspace*{2em}setptr ult\_buf, prg\_name\
\hspace*{2em}setword ult\_addr, 0\
\hspace*{2em}jsr ultimate\_load\
\hspace*{2em}jsr ultimate\_last\_end\
\hspace*{2em}sta load\_end\
\hspace*{2em}stx load\_end + 1\
\
\hspace*{2em}setword ult\_addr, \$C000\
\hspace*{2em}jsr ultimate\_load\
\
\hspace*{2em}setptr ult\_buf, raw\_name\
\hspace*{2em}setword ult\_addr, \$C000\
\hspace*{2em}setword ult\_max, 2048\
\hspace*{2em}jsr ultimate\_bload\
\
\hspace*{2em}setptr ult\_buf, save\_name\
\hspace*{2em}setword ult\_addr, \$0400\
\hspace*{2em}setword ult\_max, 1000\
\hspace*{2em}jsr ultimate\_save\
\
prg\_name:  .byte "game.prg", 0\
raw\_name:  .byte "font.bin", 0\
save\_name: .byte "screen.bin", 0\
load\_end:  .res 2
\end{listingbox}

\begin{outputbox}
load_end = address after loaded PRG data\\
FONT.BIN loads at \$C000; SCREEN.BIN receives 1000 bytes.
\end{outputbox}

An \cmd{ult\_addr} of zero makes load use the PRG header; \cmd{\$C000} shows
the explicit-address form. \cmd{ultimate\_last\_end} returns the previous load
end in \cmd{A}/\cmd{X}. Bload consumes no header and obeys \cmd{ult\_max}; save
uses \cmd{ult\_max} as its length.

\evidence{\cmd{docs/asm-abi.md}; \cmd{src/uci/file.s}}

\section{Mount A Disk Image}

The drive is named by the number it answers to on the IEC serial bus, 8 or 9,
rather than by a slot. \cmd{ULTIMATE\_DRIVE\_LAST} is zero and asks for
whichever drive was mounted into last, or drive~A when nothing has been.

\begin{listingbox}
\hspace*{2em}setptr ult\_buf, image\_name\
\hspace*{2em}lda \#8\
\hspace*{2em}jsr ultimate\_mount\
\hspace*{2em}cmp \#ULTIMATE\_OK\
\hspace*{2em}bne failed\
\
\hspace*{2em}lda \#8\
\hspace*{2em}jsr ultimate\_swap\
\
\hspace*{2em}lda \#8\
\hspace*{2em}jsr ultimate\_unmount\
\
image\_name: .byte "games/pitfall.d64", 0
\end{listingbox}

\begin{outputbox}
PITFALL.D64 is mounted, swapped and unmounted; A = ULTIMATE_OK each time.
\end{outputbox}

The device number goes in \cmd{A} and is reloaded before each call, because
every entry point returns its result there. \cmd{ultimate\_mount} takes the
image name through \cmd{ult\_buf} and measures it itself.

The firmware picks the image type from the extension: \cmd{.D64}, \cmd{.D71}
and \cmd{.D81} hold the sectors of a 1541, 1571 and 1581 disk, and \cmd{.G64}
and \cmd{.G71} hold raw track data. Mounting into a drive that is switched off
returns \cmd{ULTIMATE\_ERR\_DEVICE}.

\evidence{\cmd{bindings/asm/ultimate.inc}; \cmd{src/uci/disk.s};
\cmd{docs/uci.md}}

\Needspace{30\baselineskip}
\section{Control The Machine And Its Drives}

\begin{listingbox}
\hspace*{2em}lda \#<drives\
\hspace*{2em}ldx \#>drives\
\hspace*{2em}jsr ultimate\_drive\_info\
\hspace*{2em}cmp \#ULTIMATE\_OK\
\hspace*{2em}bne failed\
\hspace*{2em}lda drives + CTRL\_DRVINFO\_COUNT\
\hspace*{2em}beq no\_drives\
\hspace*{2em}lda drives + CTRL\_DRVINFO\_FIRST + 1\
\hspace*{2em}sta first\_device\
\
\hspace*{2em}lda \#ULTIMATE\_DRIVE\_A\
\hspace*{2em}ldx \#1\
\hspace*{2em}jsr ultimate\_drive\_enable\
\
\hspace*{2em}lda \#ULTIMATE\_DRIVE\_A\
\hspace*{2em}jsr ultimate\_drive\_power\
\hspace*{2em}lda ult\_stage\
\hspace*{2em}beq drive\_is\_off\
\
drives:       .res CTRL\_DRVINFO\_BYTES\
first\_device: .res 1
\end{listingbox}

\begin{outputbox}
drives holds a count and one record per drive; ult_stage = 1 when it runs.
\end{outputbox}

\cmd{ultimate\_drive\_info} takes the block pointer in \cmd{A}/\cmd{X}. The
block is a count at \cmd{CTRL\_DRVINFO\_COUNT}, then
\cmd{CTRL\_DRVINFO\_RECORD} bytes per drive from \cmd{CTRL\_DRVINFO\_FIRST}:
type, then the IEC device number, then 1 while the drive is running. On any
failure the count is set to zero, so a caller that skips the result check reads
no records.

\cmd{ultimate\_drive\_enable} takes the drive in \cmd{A} and the state in
\cmd{X}: zero switches it off, and any non-zero value switches it on.
\cmd{ultimate\_drive\_power} takes the drive in \cmd{A} and reports the state
in \cmd{ult\_stage}, which is 1 while the drive runs.

\begin{note}
Two numbering schemes meet here. \cmd{ultimate\_drive\_enable} and
\cmd{ultimate\_drive\_power} take \cmd{ULTIMATE\_DRIVE\_A} or
\cmd{ULTIMATE\_DRIVE\_B}, which name a physical drive.
\cmd{ultimate\_mount} takes the device number the drive answers to on the IEC
bus. \cmd{ultimate\_drive\_info} connects the two: each record carries both.
\end{note}

\begin{warning}
\cmd{ultimate\_drive\_power} writes the flag to \cmd{ult\_stage} and uses
\cmd{ult\_stage + 1} to \cmd{ult\_stage + 3} as scratch. Read the flag before
the next SDK call.
\end{warning}

\subsection*{Reset the machine, or open its menu}

\begin{listingbox}
\hspace*{2em}jsr ultimate\_freeze\
\
\hspace*{2em}jsr ultimate\_reboot\
\hspace*{2em}; nothing here runs
\end{listingbox}

\begin{outputbox}
freeze returns when the menu is left; reboot does not return.
\end{outputbox}

\cmd{ultimate\_reboot} resets the C64, so the program calling it stops existing
part way through the call. \cmd{ultimate\_freeze} stops the C64 and shows the
Ultimate's menu, and returns when whoever is at the keyboard leaves it.

\subsection*{Read the RAM disk layout}

\begin{listingbox}
\hspace*{2em}lda \#<ramdisk\
\hspace*{2em}ldx \#>ramdisk\
\hspace*{2em}jsr ultimate\_ramdisk\_info\
\
ramdisk: .res CTRL\_RAMDISK\_BYTES
\end{listingbox}

\begin{outputbox}
ramdisk holds CTRL_RAMDISK_DRIVES two-byte records.
\end{outputbox}

Each record is a drive type code and a size in 64K units. A type of zero is a
slot with nothing in it. These are the RAM disks GEOS MegaPatch uses.

\evidence{\cmd{bindings/asm/ultimate.inc}; \cmd{src/uci/control.s};
\cmd{docs/uci.md}}

\Needspace{30\baselineskip}
\section{Read And Set The Clock}

\begin{listingbox}
\hspace*{2em}setptr ult\_buf, time\_text\
\hspace*{2em}setword ult\_buflen, ULTIMATE\_TIME\_BUFFER\
\hspace*{2em}setword ult\_outlen, 0\
\hspace*{2em}lda \#ULTIMATE\_TIME\_PLAIN\
\hspace*{2em}jsr ultimate\_get\_time\
\
time\_text: .res ULTIMATE\_TIME\_BUFFER
\end{listingbox}

\begin{outputbox}
time_text = "2026/08/22 14:30:00"
\end{outputbox}

\cmd{A} is the format. \cmd{ULTIMATE\_TIME\_PLAIN} asks for the date and time
and \cmd{ULTIMATE\_TIME\_WEEKDAY} puts the weekday in front of it.
\cmd{ULTIMATE\_TIME\_BUFFER} is a buffer size that fits either form. This is
one of the few entry points that reads \cmd{ult\_buflen} as an input rather
than computing it. A buffer too small still returns \cmd{ULTIMATE\_OK}, with
the part that fitted.

\subsection*{Set it}

\begin{listingbox}
\hspace*{2em}lda \#126\hspace{2em}; 2026, as the year less 1900\
\hspace*{2em}sta ult\_stage + 0\
\hspace*{2em}lda \#8\hspace{3em}; month\
\hspace*{2em}sta ult\_stage + 1\
\hspace*{2em}lda \#22\hspace{2em}; day\
\hspace*{2em}sta ult\_stage + 2\
\hspace*{2em}lda \#14\hspace{2em}; hour\
\hspace*{2em}sta ult\_stage + 3\
\hspace*{2em}lda \#30\hspace{2em}; minute\
\hspace*{2em}sta ult\_stage + 4\
\hspace*{2em}lda \#0\hspace{3em}; second\
\hspace*{2em}sta ult\_stage + 5\
\hspace*{2em}jsr ultimate\_set\_time
\end{listingbox}

\begin{outputbox}
The clock is set to 2026/08/22 14:30:00; A = ULTIMATE_OK.
\end{outputbox}

The six bytes go in \cmd{ult\_stage} in that order. \cmd{ult\_stagelen} is not
a caller input here; the command always sends six bytes.

\begin{warning}
\cmd{ult\_stage} is shared. \cmd{ultimate\_mount}, \cmd{ultimate\_unmount},
\cmd{ultimate\_swap}, \cmd{ultimate\_get\_time},
\cmd{ultimate\_drive\_info} and \cmd{ultimate\_http\_body} all overwrite
\cmd{ult\_stage + 0}, and \cmd{ultimate\_drive\_power} and the keyed body calls
overwrite more of it. Fill it immediately before \cmd{ultimate\_set\_time}.
\end{warning}

\evidence{\cmd{bindings/asm/ultimate.inc}; \cmd{src/uci/clock.s};
\cmd{docs/uci.md}}

\section{Move Bytes Through The REU}

\begin{listingbox}
\hspace*{2em}jsr ultimate\_reu\_available\
\hspace*{2em}beq no\_reu\
\hspace*{2em}jsr ultimate\_reu\_size\
\hspace*{2em}sta reu\_banks\
\hspace*{2em}stx reu\_banks + 1\
\
\hspace*{2em}setword ult\_addr, \$0400\
\hspace*{2em}lda \#0\
\hspace*{2em}sta ult\_reu\
\hspace*{2em}sta ult\_reu + 1\
\hspace*{2em}sta ult\_reu + 2\
\hspace*{2em}sta ult\_reu + 3\
\hspace*{2em}setword ult\_reulen, 1000\
\hspace*{2em}lda \#0\
\hspace*{2em}sta ult\_reulen + 2\
\hspace*{2em}sta ult\_reulen + 3\
\hspace*{2em}jsr ultimate\_reu\_stash\
\hspace*{2em}jsr ultimate\_reu\_fetch\
\
reu\_banks: .res 2
\end{listingbox}

\begin{outputbox}
reu_banks = installed REU size in 64K banks\\
1000 bytes stashed and fetched; A = ULTIMATE_OK.
\end{outputbox}

Size returns 64K banks in \cmd{A}/\cmd{X}. Stash copies C64 RAM to the REU;
fetch copies it back. Both read the C64 address, 32-bit REU address and 32-bit
length from the shared variables.

\subsection*{Move the current file directly to and from the REU}

Both transfers below use REU address zero and a 65536-byte length:

\begin{listingbox}
\hspace*{2em}lda \#0\
\hspace*{2em}sta ult\_reu\
\hspace*{2em}sta ult\_reu + 1\
\hspace*{2em}sta ult\_reu + 2\
\hspace*{2em}sta ult\_reu + 3\
\hspace*{2em}sta ult\_reulen\
\hspace*{2em}sta ult\_reulen + 1\
\hspace*{2em}lda \#1\
\hspace*{2em}sta ult\_reulen + 2\
\hspace*{2em}lda \#0\
\hspace*{2em}sta ult\_reulen + 3
\end{listingbox}

\begin{outputbox}
ult_reu = \$00000000; ult_reulen = 65536.
\end{outputbox}

\Needspace{25\baselineskip}
\subsection*{File to REU}

\begin{listingbox}
\hspace*{2em}; Open ASSETS.BIN first; it becomes the current DOS file.\
\hspace*{2em}setptr ult\_buf, asset\_name\
\hspace*{2em}lda \#DOS\_FA\_READ\
\hspace*{2em}jsr ultimate\_open\
\hspace*{2em}cmp \#ULTIMATE\_OK\
\hspace*{2em}bne failed\
\hspace*{2em}lda \#0\
\hspace*{2em}sta ult\_num\
\hspace*{2em}sta ult\_num + 1\
\hspace*{2em}sta ult\_num + 2\
\hspace*{2em}sta ult\_num + 3\
\hspace*{2em}jsr ultimate\_seek\
\hspace*{2em}jsr ultimate\_reu\_load\
\hspace*{2em}jsr ultimate\_close\
\
asset\_name: .byte "assets.bin", 0
\end{listingbox}

\begin{outputbox}
ASSETS.BIN to REU \$000000: 65536 bytes
\end{outputbox}

\Needspace{18\baselineskip}
\subsection*{REU to file}

\begin{listingbox}
\hspace*{2em}; Open a writable current file before the transfer.\
\hspace*{2em}setptr ult\_buf, capture\_name\
\hspace*{2em}lda \#DOS\_FA\_CREATE\_ALWAYS | DOS\_FA\_WRITE\
\hspace*{2em}jsr ultimate\_open\
\hspace*{2em}cmp \#ULTIMATE\_OK\
\hspace*{2em}bne failed\
\hspace*{2em}jsr ultimate\_reu\_save\
\hspace*{2em}jsr ultimate\_close\
\
capture\_name: .byte "capture.bin", 0
\end{listingbox}

\begin{outputbox}
REU \$000000 to CAPTURE.BIN: 65536 bytes
\end{outputbox}

The REU address is zero and the length bytes encode 65536. Load and save take
no filename: each operates on the firmware's current open file established by
the preceding \cmd{ultimate\_open}.

\evidence{\cmd{docs/asm-abi.md}; \cmd{src/uci/reu.s}}

\Needspace{32\baselineskip}
\section{Inspect Network Interfaces}

\begin{listingbox}
\hspace*{2em}jsr ultimate\_net\_ifcount\
\hspace*{2em}cmp \#ULTIMATE\_OK\
\hspace*{2em}bne failed\
\hspace*{2em}lda ult\_iface\
\hspace*{2em}beq no\_interfaces\
\
\hspace*{2em}lda \#0\
\hspace*{2em}sta ult\_iface\
\hspace*{2em}setptr ult\_buf, mac\
\hspace*{2em}jsr ultimate\_net\_macaddr\
\hspace*{2em}setptr ult\_buf, ipconfig\
\hspace*{2em}jsr ultimate\_net\_ipconfig\
\
mac:      .res 6\
ipconfig: .res 12
\end{listingbox}

\begin{outputbox}
ult_iface = interface count; mac receives 6 bytes; ipconfig receives 12.
\end{outputbox}

The count is returned in \cmd{ult\_iface}; the address calls take an interface
index there and place fixed-size replies at \cmd{ult\_buf}.

\subsection*{Change an interface's address}

\begin{listingbox}
\hspace*{2em}lda \#0\
\hspace*{2em}sta ult\_iface\
\hspace*{2em}setptr ult\_buf, ipconfig\
\hspace*{2em}jsr ultimate\_net\_ipconfig\
\hspace*{2em}cmp \#ULTIMATE\_OK\
\hspace*{2em}bne failed\
\
\hspace*{2em}lda \#60\
\hspace*{2em}sta ipconfig + 3\
\hspace*{2em}jsr ultimate\_net\_setip
\end{listingbox}

\begin{outputbox}
The interface keeps its netmask and gateway and answers on .60 instead.
\end{outputbox}

\cmd{ultimate\_net\_setip} sends back the same twelve bytes
\cmd{ultimate\_net\_ipconfig} returned: four of address, then four of netmask,
then four of gateway. It reads \cmd{ult\_iface} and \cmd{ult\_buf}, and always
sends twelve bytes, so \cmd{ult\_buflen} is not an input.

This changes the running configuration only; nothing is written to the
Ultimate's stored settings.

\section{Use TCP And UDP Sockets}

\begin{listingbox}
\hspace*{2em}setptr ult\_buf, host\
\hspace*{2em}setword ult\_port, 6502\
\hspace*{2em}jsr ultimate\_net\_connect\
\hspace*{2em}cmp \#ULTIMATE\_OK\
\hspace*{2em}bne failed\
\
\hspace*{2em}setptr ult\_buf, message\
\hspace*{2em}setword ult\_buflen, message\_end - message\
\hspace*{2em}jsr ultimate\_net\_write\
\hspace*{2em}; ult\_socklen is now the number accepted.\
\
host:        .byte "192.168.1.50", 0\
message:     .byte "hi", 13\
message\_end:
\end{listingbox}

\begin{outputbox}
TCP handle is in ult_sock; ult_socklen = 3 bytes accepted.
\end{outputbox}

\Needspace{25\baselineskip}
\subsection*{Read and close the TCP socket}

\begin{listingbox}
poll:\
\hspace*{2em}setptr ult\_buf, socket\_buffer\
\hspace*{2em}setword ult\_socklen, 128\
\hspace*{2em}jsr ultimate\_net\_read\
\hspace*{2em}cmp \#ULTIMATE\_END\
\hspace*{2em}beq peer\_closed\
\hspace*{2em}cmp \#ULTIMATE\_OK\
\hspace*{2em}bne close\_socket\
\hspace*{2em}lda ult\_socklen\
\hspace*{2em}ora ult\_socklen + 1\
\hspace*{2em}beq poll\
\
close\_socket:\
\hspace*{2em}jsr ultimate\_net\_close\
\
socket\_buffer: .res 128
\end{listingbox}

\begin{outputbox}
Reply bytes are in socket_buffer; close returns ULTIMATE_OK.
\end{outputbox}

\Needspace{14\baselineskip}
\subsection*{UDP}

\begin{listingbox}
\hspace*{2em}setptr ult\_buf, host\
\hspace*{2em}setword ult\_port, 6502\
\hspace*{2em}jsr ultimate\_net\_udp\
\hspace*{2em}cmp \#ULTIMATE\_OK\
\hspace*{2em}bne failed\
\hspace*{2em}jsr ultimate\_net\_close
\end{listingbox}

\begin{outputbox}
UDP handle opens and closes; A = ULTIMATE_OK.
\end{outputbox}

Connect and UDP return the handle in \cmd{ult\_sock}. Write takes its length
from \cmd{ult\_buflen} and reports accepted bytes in \cmd{ult\_socklen}. Read
takes capacity and returns stored bytes through \cmd{ult\_socklen}. A zero-byte
success means poll again. \cmd{ULTIMATE\_END} means the handle is already gone;
do not close that ended handle.

\evidence{\cmd{docs/asm-abi.md}; \cmd{src/uci/net.s}}

\section{Use The HTTP Client}

\begin{listingbox}
\hspace*{2em}setptr ult\_url, url\
\hspace*{2em}setptr ult\_buf, http\_reply\
\hspace*{2em}setword ult\_buflen, 512\
\hspace*{2em}jsr ultimate\_http\_get\
\hspace*{2em}; ult\_httplen is the number stored.\
\
url:           .byte "http://192.168.1.50/", 0\
accept\_header: .byte "accept: text/plain", 0\
http\_reply:    .res 512
\end{listingbox}

\begin{outputbox}
One-call GET stores response bytes; A = ULTIMATE_OK.
\end{outputbox}

\Needspace{26\baselineskip}
\subsection*{Multi-step GET without a body}

\begin{listingbox}
\hspace*{2em}setptr ult\_url, url\
\hspace*{2em}lda \#HTTP\_VERB\_GET\
\hspace*{2em}jsr ultimate\_http\_open\
\hspace*{2em}cmp \#ULTIMATE\_OK\
\hspace*{2em}bne failed\
\
\hspace*{2em}setptr ult\_url, accept\_header\
\hspace*{2em}jsr ultimate\_http\_header\
\hspace*{2em}lda \#HTTP\_BODY\_NONE\
\hspace*{2em}sta ult\_body\
\hspace*{2em}setptr ult\_buf, http\_reply\
\hspace*{2em}setword ult\_buflen, 512\
\hspace*{2em}jsr ultimate\_http\_exchange\
\hspace*{2em}jsr ultimate\_http\_close\
\
\end{listingbox}

\begin{outputbox}
GET without a body stores response bytes; close returns ULTIMATE_OK.
\end{outputbox}

\Needspace{22\baselineskip}
The one-call GET creates and releases its own slot. The multi-step form stores
the request handle in \cmd{ult\_http}, and header, exchange and close all read
it there. \cmd{ult\_body} holds the body handle, and
\cmd{HTTP\_BODY\_NONE} in it means no body.
\cmd{ultimate\_http\_free\_all} recovers slots abandoned by an earlier
program.

The next section creates a body and posts it.

\evidence{\cmd{docs/asm-abi.md}; \cmd{src/uci/http.s}}

\Needspace{32\baselineskip}
\section{Build A Request Body}

The body is built inside the Ultimate: create a slot, add to it one call at a
time, then run the exchange. The handle lives in \cmd{ult\_body}, where
\cmd{ultimate\_http\_exchange} already looks for it, so nothing has to carry it
around.

\begin{listingbox}
\hspace*{2em}lda \#HTTP\_BODY\_JSON\_OBJECT\
\hspace*{2em}jsr ultimate\_http\_body\
\hspace*{2em}cmp \#ULTIMATE\_OK\
\hspace*{2em}bne failed\
\
\hspace*{2em}setptr ult\_url, key\_name\
\hspace*{2em}setptr ult\_buf, value\_ada\
\hspace*{2em}jsr ultimate\_http\_body\_string\
\
\hspace*{2em}setptr ult\_url, key\_score\
\hspace*{2em}setword ult\_val, 4200\
\hspace*{2em}lda \#0\
\hspace*{2em}sta ult\_val + 2\
\hspace*{2em}sta ult\_val + 3\
\hspace*{2em}jsr ultimate\_http\_body\_int\
\
key\_name:  .byte "name", 0\
key\_score: .byte "score", 0\
value\_ada: .byte "ada", 0
\end{listingbox}

\begin{outputbox}
The body holds \{"NAME":"ADA","SCORE":4200\}.
\end{outputbox}

\cmd{ult\_url} is the key for every keyed call, and it is the same variable
\cmd{ultimate\_http\_open} uses for a URL. \cmd{ult\_buf} is a string value.
\cmd{ult\_val} is four bytes, low byte first, and all four are sent, because
the firmware sign-extends from the last byte it receives.

\cmd{ultimate\_http\_body\_bool} reads \cmd{ult\_val + 0} only, and any
non-zero value is true.

\subsection*{Send it}

\begin{listingbox}
\hspace*{2em}setptr ult\_url, url\
\hspace*{2em}lda \#HTTP\_VERB\_POST\
\hspace*{2em}jsr ultimate\_http\_open\
\hspace*{2em}cmp \#ULTIMATE\_OK\
\hspace*{2em}bne failed\
\
\hspace*{2em}setptr ult\_buf, http\_reply\
\hspace*{2em}setword ult\_buflen, 512\
\hspace*{2em}jsr ultimate\_http\_exchange\
\hspace*{2em}jsr ultimate\_http\_close\
\hspace*{2em}jsr ultimate\_http\_body\_free
\end{listingbox}

\begin{outputbox}
The response is in http_reply; ult_httplen is the number stored.
\end{outputbox}

Set \cmd{ult\_buf} back to the reply buffer before the exchange. The keyed body
calls used it for the last string value, and the exchange uses it for the
reply.

\subsection*{Nest an object or an array}

\begin{listingbox}
\hspace*{2em}setptr ult\_url, key\_stats\
\hspace*{2em}jsr ultimate\_http\_body\_object\
\hspace*{2em}; keys added here go inside "stats"\
\hspace*{2em}jsr ultimate\_http\_body\_up\
\
\hspace*{2em}setptr ult\_url, key\_runs\
\hspace*{2em}jsr ultimate\_http\_body\_array\
\hspace*{2em}; inside an array the key is ignored\
\hspace*{2em}jsr ultimate\_http\_body\_up
\end{listingbox}

\begin{outputbox}
\{"STATS":\{...\},"RUNS":[...]\}
\end{outputbox}

Adding an object or an array enters it. Every key added afterwards goes inside,
until \cmd{ultimate\_http\_body\_up} steps back out to the parent.

\subsection*{Send raw bytes instead}

\begin{listingbox}
\hspace*{2em}lda \#HTTP\_BODY\_BINARY\
\hspace*{2em}jsr ultimate\_http\_body\
\
\hspace*{2em}setptr ult\_buf, block\
\hspace*{2em}setword ult\_buflen, block\_end - block\
\hspace*{2em}jsr ultimate\_http\_body\_binary\
\
\hspace*{2em}jsr ultimate\_http\_body\_clear\
\hspace*{2em}jsr ultimate\_http\_body\_free\
\
block:     .byte 1, 2, 3, 4\
block\_end:
\end{listingbox}

\begin{outputbox}
The body holds four bytes, then none; the slot is returned.
\end{outputbox}

\cmd{ultimate\_http\_body\_binary} is the other entry that reads
\cmd{ult\_buflen} as an input. Successive calls append.
\cmd{ultimate\_http\_body\_clear} empties a body and keeps its slot and format;
\cmd{ultimate\_http\_body\_free} returns the slot.

A \cmd{HTTP\_BODY\_BINARY} body takes \cmd{ultimate\_http\_body\_binary} and
nothing else; the JSON and form formats take everything else. A key is staged
through \cmd{ult\_stage}, so it is limited to \cmd{ULTIMATE\_HTTP\_KEY\_MAX}
bytes, which is 34. A longer key returns
\cmd{ULTIMATE\_ERR\_INVALID\_ARGUMENT} and never reaches the wire.

\begin{warning}
The firmware has sixteen slots for the whole machine, and a program that
crashes returns none of them. Free every body and every request.
\cmd{ultimate\_http\_free\_all} takes back headers and bodies together.
\end{warning}

\evidence{\cmd{bindings/asm/ultimate.inc}; \cmd{src/uci/httpbody.s}}

\section{Probe And Configure The Raw Transport}

\begin{listingbox}
\hspace*{2em}jsr uci\_present\
\hspace*{2em}beq no\_interface\
\hspace*{2em}jsr uci\_ident\
\hspace*{2em}sta raw\_ident\
\
\hspace*{2em}jsr uci\_get\_timeout\_a\
\hspace*{2em}pha\
\hspace*{2em}lda \#UCI\_TIMEOUT\_FOREVER\
\hspace*{2em}jsr uci\_set\_timeout\_a\
\hspace*{2em}jsr run\_slow\_command\
\hspace*{2em}tax\
\hspace*{2em}pla\
\hspace*{2em}jsr uci\_set\_timeout\_a\
\hspace*{2em}txa\
\
\hspace*{2em}jsr uci\_last\_code\
\hspace*{2em}sta device\_code\
\hspace*{2em}stx device\_code + 1\
\hspace*{2em}jsr uci\_abort
\end{listingbox}

\begin{outputbox}
uci_present returns nonzero; raw_ident receives the interface ID.\\
device_code receives the last target status; abort returns ULTIMATE_OK.
\end{outputbox}

Present and ident are direct register reads. Timeout zero waits forever; this
pattern restores the previous budget while preserving the command result.
Last code returns the raw target status in \cmd{A}/\cmd{X}. Abort releases an
unfinished exchange and is safe while idle.

\section{Raw Commands: Optional Arguments}

\begin{listingbox}
\hspace*{2em}jsr uci\_req\_clear\
\hspace*{2em}lda \#UCI\_TARGET\_DOS1\
\hspace*{2em}sta uci\_req + UCI\_REQ\_TARGET\
\hspace*{2em}lda \#UCI\_CMD\_IDENTIFY\
\hspace*{2em}sta uci\_req + UCI\_REQ\_COMMAND\
\hspace*{2em}setptr uci\_req + UCI\_REQ\_DATA, raw\_reply\
\hspace*{2em}setword uci\_req + UCI\_REQ\_DATAMAX, 64\
\hspace*{2em}jsr uci\_exec\_block\
\hspace*{2em}jsr uci\_req\_clear\
\hspace*{2em}lda \#UCI\_TARGET\_DOS1\
\hspace*{2em}sta uci\_req + UCI\_REQ\_TARGET\
\hspace*{2em}lda \#DOS\_CMD\_GET\_TIME\
\hspace*{2em}sta uci\_req + UCI\_REQ\_COMMAND\
\hspace*{2em}setptr uci\_req + UCI\_REQ\_ARGS, time\_format\
\hspace*{2em}setword uci\_req + UCI\_REQ\_ARGLEN, 1\
\hspace*{2em}jsr uci\_exec\_block\
\hspace*{2em}jsr uci\_req\_clear\
\hspace*{2em}lda \#UCI\_TARGET\_CONTROL\
\hspace*{2em}sta uci\_req + UCI\_REQ\_TARGET\
\hspace*{2em}lda \#CTRL\_CMD\_GET\_HWINFO\
\hspace*{2em}sta uci\_req + UCI\_REQ\_COMMAND\
\hspace*{2em}setptr uci\_req + UCI\_REQ\_DATA, raw\_reply\
\hspace*{2em}setword uci\_req + UCI\_REQ\_DATAMAX, 64\
\hspace*{2em}jsr uci\_exec\_block\
\hspace*{2em}setptr uci\_req + UCI\_REQ\_ARGS, hwinfo\_model\
\hspace*{2em}setword uci\_req + UCI\_REQ\_ARGLEN, 1\
\hspace*{2em}jsr uci\_exec\_block\
time\_format: .byte 0\
hwinfo\_model: .byte CTRL\_HWINFO\_MODEL\
raw\_reply:   .res 64
\end{listingbox}

\begin{outputbox}
IDENTIFY returns the target name.\\
GET_TIME returns the formatted time.\\
Both GET_HWINFO forms return the model name.
\end{outputbox}

Clearing the request leaves every optional span empty. The time command then
supplies its optional format byte. \cmd{GET\_HWINFO} is sent without and then
with its model selector, which is what an omitted optional argument means on
the wire. Set other spans only when needed.

\subsection*{Use a caller-owned request block}

\begin{listingbox}
\hspace*{2em}lda \#<my\_request\
\hspace*{2em}ldx \#>my\_request\
\hspace*{2em}jsr uci\_exec\
\
my\_request: .res UCI\_REQ\_SIZE
\end{listingbox}

\begin{outputbox}
A = request result; lengths and retained bytes are written into my_request.
\end{outputbox}

\cmd{uci\_exec} takes the request pointer directly. Use it instead of the
global block when the program owns more than one request.

\Needspace{30\baselineskip}
\section{Read A Raw Multi-block Reply}

\begin{listingbox}
\hspace*{2em}lda \#<my\_request\
\hspace*{2em}ldx \#>my\_request\
\hspace*{2em}jsr uci\_exec\_first\
\hspace*{2em}cmp \#ULTIMATE\_OK\
\hspace*{2em}bne failed\
\
next\_block:\
\hspace*{2em}; use this block's UCI\_REQ\_DATALEN bytes\
\hspace*{2em}lda uci\_more\
\hspace*{2em}beq all\_blocks\
\hspace*{2em}jsr uci\_exec\_next\
\hspace*{2em}cmp \#ULTIMATE\_OK\
\hspace*{2em}beq next\_block
\end{listingbox}

\begin{outputbox}
Each iteration receives one block; uci_more = 0 after the final block.
\end{outputbox}

First submits the request and returns one block. \cmd{uci\_more} is the boolean
state byte that says whether another follows; next releases the current block
and receives it. Issue no unrelated command until the final block or call
\cmd{uci\_abort} when stopping early.

\evidence{\cmd{docs/asm-abi.md}; \cmd{src/uci/uci\_core.s}}

% Chapter 5 - the standalone blob.

\guidechapter{OTHER TOOLCHAINS}{%
  \item Building and identifying the blob
  \item The stable jump table
  \item The parameter block
  \item Ultimate Audio and software-vsprite extension entries
  \item Calling it from BASIC
  \item Choosing a different base address
}

\section{Build And Identify The Blob}

The blob is the SDK linked as one binary with a table of jump instructions at
a fixed address. Any assembler or compiler can call it, because calling it
needs no linker and no object format: you load the binary and \cmd{jsr} into
the table.

\cmd{make blob} produces \cmd{bindings/blob/build/ultimate-7000.bin}. Load it
at \cmd{\$7000}. The first three bytes are the PETSCII signature \cmd{UCI};
the fourth is the version of the ABI, the calling interface.

\begin{note}
The blob has to sit below \cmd{\$A000}, because \cmd{\$A000--\$BFFF} is BASIC
ROM. A program can reach every byte of a blob at \cmd{\$7000} with the machine
exactly as it found it. One that ran past \cmd{\$A000} would have to switch
BASIC ROM out and back around every call.
\end{note}

\begin{warning}
\cmd{\$7000} is inside BASIC's memory, not above it. BASIC owns
\cmd{\$0801--\$9FFF}: the program area grows up from \cmd{\$0801} and strings
and arrays grow down from \cmd{\$9FFF}. A program that has replaced BASIC owns
that memory instead and needs to do nothing. A program that leaves BASIC running
must lower the top of BASIC memory to the base address before it loads the blob,
or BASIC will write over it. \emph{Call The Blob From BASIC}, later in this
chapter, shows the two \bas{poke}s and the \bas{clr} that do it.
\end{warning}

This example checks the default build before calling it:

\begin{listingbox}
\hspace*{2em}lda \$7000\\
\hspace*{2em}cmp \#\$D5\hspace{2em}; PETSCII 'U'\\
\hspace*{2em}bne no\_sdk\\
\hspace*{2em}lda \$7001\\
\hspace*{2em}cmp \#\$C3\hspace{2em}; PETSCII 'C'\\
\hspace*{2em}bne no\_sdk\\
\hspace*{2em}lda \$7002\\
\hspace*{2em}cmp \#\$C9\hspace{2em}; PETSCII 'I'\\
\hspace*{2em}bne no\_sdk\\
\hspace*{2em}lda \$7003\\
\hspace*{2em}cmp \#1\hspace{3em}; ABI version\\
\hspace*{2em}bne wrong\_version
\end{listingbox}

\begin{resultbox}
All four comparisons match; execution continues with ABI version 1.
\end{resultbox}

The version byte changes only when an existing entry changes meaning. New
entries are appended, so a program written against version~1 keeps working
when the table grows.

\evidence{\cmd{bindings/blob/Makefile}; \cmd{bindings/blob/blob.s};
\cmd{tests/emulator/blob.suite}}

\section{Call The Stable Jump Table}

Entries begin at base plus \cmd{\$04}. Each entry is a three-byte \cmd{jmp};
new entries are appended so published offsets do not move.

\begin{listingbox}
\hspace*{2em}jsr \$7004\hspace{2em}; +\$04  uci\_init\\
\hspace*{2em}jsr \$701C\hspace{2em}; +\$1C  ultimate\_init
\end{listingbox}

\begin{resultbox}
A = ULTIMATE_OK after each initialization call.
\end{resultbox}

These are the two initialization entries. Use \cmd{ultimate\_init} for
application code; the lower-level \cmd{uci\_init} entry is exposed for
transport users. Both return a result in \cmd{A}.

\subsection*{The two kinds of entry}

The table has two halves, and the difference decides how you read a result.

Entries from \cmd{+\$04} to \cmd{+\$4C} pass their arguments in registers, the
way the assembly chapter's entry points do, and return the result in \cmd{A}
and nowhere else.

Entries from \cmd{+\$4F} upwards take their arguments from the parameter block
described in the next section, because a caller cannot pass a filename in
\cmd{A} and \cmd{X}. Each of these also stores its result in \cmd{bp\_result}
before returning it in \cmd{A}.

That second group is what a BASIC program can drive, because \bas{sys} can
neither pass a register nor read one back. The extension tables at
\cmd{+\$2E8} and \cmd{+\$300} contain both kinds; their individual contracts
say whether they use registers or the parameter block.

The complete offset table is in \cmd{bindings/blob/README.md}. Read the
offsets there rather than counting entries: the table gained the file
information, disk image, clock, machine and HTTP body services after the
offsets printed in this chapter were published.

\evidence{\cmd{bindings/blob/blob.s}; \cmd{bindings/blob/README.md};
\cmd{tests/emulator/blob.suite}}

\Needspace{20\baselineskip}
\section{Use The Parameter Block}

Service entries that need several arguments use a 512-byte block at
base plus \cmd{\$100}. These are the first four fields:

\begin{reftable}{llp{0.48\linewidth}}
Offset & Field & Purpose \\
\midrule
\cmd{+\$00} & \cmd{bp\_result} & last service result \\
\cmd{+\$05} & \cmd{bp\_len} & length in or out \\
\cmd{+\$29} & \cmd{bp\_name} & 40-byte area for a terminated name \\
\cmd{+\$51} & \cmd{bp\_reply} & 256-byte reply area \\
\end{reftable}

Later fields extend beyond the first page while retaining their published
offsets. Consult the blob README for the complete layout rather than deriving
an address from this subset.

\evidence{\cmd{bindings/blob/blob.s}; \cmd{bindings/blob/README.md};
\cmd{tests/emulator/blob.suite}}

\Needspace{34\baselineskip}
\section{Call Audio And Vsprite Extensions}

Ultimate Audio occupies the extension entries from \cmd{+\$2E8} through
\cmd{+\$300}. The normal sequence is
\cmd{ultimate\_audio\_init}, optionally
\cmd{audio\_load\_wav}, then \cmd{audio\_configure} and
\cmd{audio\_start}. Configure, start, stop and clear use the
\cmd{ultimate\_audio\_voice} structure at
\cmd{BLOB + BLOB\_BP\_AUDIO}; the WAV entry also uses
\cmd{BLOB\_BP\_NAME} and \cmd{BLOB\_BP\_REU}. Chapter~3 defines every voice
field and the safe configure/start/stop sequence. The complete stable offset
table is in \cmd{bindings/blob/README.md}.

The software-vsprite entry is \cmd{BLOB + BLOB\_VSPRITE\_DRAW}, currently
\cmd{+\$303}. Put the descriptor address in \cmd{BLOB\_BP\_ADDR} and read the
result from \cmd{BLOB\_BP\_RESULT}:

\begin{listingbox}
BLOB = \$7000\
\
\hspace*{2em}lda \#<sprite\
\hspace*{2em}sta BLOB + BLOB\_BP\_ADDR\
\hspace*{2em}lda \#>sprite\
\hspace*{2em}sta BLOB + BLOB\_BP\_ADDR + 1\
\hspace*{2em}jsr BLOB + BLOB\_VSPRITE\_DRAW\
\hspace*{2em}lda BLOB + BLOB\_BP\_RESULT\
\hspace*{2em}bne draw\_failed\
\
sprite:\
\hspace*{2em}.word \$6000, image, mask, 0\
\hspace*{2em}.byte 12, 72, 2, 8, 0, VSPRITE\_F\_MASKED
\end{listingbox}

\begin{resultbox}
The masked image is composited into the bitmap; the parameter result is zero.
\end{resultbox}

The descriptor and flag contract is the one in \emph{Composite A Software
Vsprite} in Chapter~3. This is a parameter-block shim around the same assembly
routine, not another renderer.

\evidence{\cmd{bindings/blob/blob.s}; \cmd{bindings/blob/README.md};
\cmd{bindings/asm/uci\_protocol.inc}; \cmd{src/uci/vsprite.s}}

\section{Call The Blob From BASIC}

The page-aligned block gives BASIC a calling convention it can express with
\bas{poke}, \bas{peek} and \bas{sys}. Before any of that, BASIC has to be told
to stay out of the memory the blob occupies. Type this first, then load the
blob:

\begin{screen}
POKE 56,112 : POKE 55,0 : CLR
\end{screen}

112 is \cmd{\$70}, the high byte of the base address, and 55 and 56 are the low
and high bytes of BASIC's top-of-memory pointer at \cmd{\$37--\$38}. \bas{clr}
is what makes it take effect, because it resets the string pointer from the new
top. BASIC then keeps \cmd{\$0801--\$6FFF}, and the blob's code at \cmd{\$7000}
and its variables at \cmd{\$9F00} are both out of its reach. Loading the blob
without this leaves it where BASIC's strings and arrays will eventually be
written.

This program then checks the signature, brings the SDK up, confirms that an
Ultimate answered, and prints the size of the RAM expansion:

\begin{screen}
10 B = 28672 : P = B + 256\\
20 IF PEEK(B)<>213 OR PEEK(B+1)<>195 THEN PRINT "NO SDK":END\\
25 IF PEEK(B+2)<>201 THEN PRINT "NO SDK":END\\
30 SYS B + 28\\
40 SYS B + 82 : IF PEEK(P) THEN PRINT "NO ULTIMATE":END\\
50 SYS B + 175\\
60 PRINT "REU BANKS:"; PEEK(P+5) + PEEK(P+6)*256
\end{screen}

Line~10 sets \bas{b} to the blob's base, decimal 28672 for \cmd{\$7000}, and
\bas{p} to the parameter block a page above it. Lines~20 and~25 compare the
three signature bytes; 213, 195 and 201 are \cmd{\$D5}, \cmd{\$C3} and
\cmd{\$C9}.

Line~30 calls \cmd{ultimate\_init} at offset \cmd{\$1C}. That is a register
entry, so it returns its result in \cmd{A} and writes nothing to the block,
and \bas{sys} cannot read \cmd{A}. Line~40 therefore calls \cmd{getpath} at
offset \cmd{\$52}, which is a parameter-block entry: it asks the Ultimate for
the current directory and leaves the result in \cmd{bp\_result}, so
\bas{peek(p)} is the first result this program can test. A non-zero value
there means no Ultimate answered.

Line~50 calls \cmd{reu\_size} at offset \cmd{\$AF}, which puts the 16-bit bank
count in \cmd{bp\_len}. Line~60 reads it back as two bytes, low byte first, and
prints it. A machine with a 2MB expansion prints 32; one with no expansion
prints 0.

\begin{warning}
Do not check \bas{peek(p)} after a register entry such as \cmd{\$1C}. Those
entries never write \cmd{bp\_result}, so the byte still holds whatever the
previous call left there, and in a freshly loaded blob that is zero. The test
would pass whether the call worked or not.
\end{warning}

\evidence{\cmd{bindings/blob/blob.s}; \cmd{bindings/blob/README.md};
\cmd{tests/emulator/blob.suite}}

\section{Choose A Different Base Address}

Prefer building at the final base: \cmd{make -C bindings/blob BASE=6000} links
a \cmd{\$6000} image directly, and \cmd{VARS} and \cmd{ZP} move the SDK's RAM
the same way.

When the address is known only at run time, the build also emits a relocation
table beside the binary for \cmd{reloc.s} to apply. The relocated emulator
suite copies the default blob to \cmd{\$6000}, relocates the copy, calls it,
and then compares it byte for byte against an image the linker built for
\cmd{\$6000} directly.

\evidence{\cmd{bindings/blob/Makefile}; \cmd{bindings/blob/reloc.s};
\cmd{tests/emulator/blob-relocated.suite.in};
\cmd{tests/emulator/relocharness.s}}

% Chapter 6 - the services.

\guidechapter{THE SERVICES}{%
  \item Files and directories
  \item Loading and saving
  \item The RAM expansion
  \item Processor speed
  \item The screen palette
  \item Network sockets
  \item Fetching over HTTP
}

\section{Files And Directories}

The Ultimate's own filesystem, reached by name. Open a file, read or write it,
seek in it, close it --- the operations you would expect, with two things worth
knowing.

A directory walk is \textbf{one live exchange}: \cmd{ultimate\_opendir()} then
\cmd{ultimate\_readdir()} until it answers \cmd{ULTIMATE\_END}, with no other
command in between. And a filename goes on the wire exactly as you gave it: the
SDK converts nothing, because uppercase PETSCII and uppercase ASCII are the
same bytes and the filesystem does not care about case.

\section{Loading And Saving}

\cmd{ultimate\_load()} takes the firmware's fast path when the machine has one
--- the bytes go straight into C64 memory without passing through the command
interface at all --- and falls back to the ordinary route when it does not. You
do not choose; the SDK does, and tells you where the file ended.

\cmd{ultimate\_bload()} is the one for data: no header consumed, nothing
interpreted, and not one byte past the limit you set.

\section{The RAM Expansion}

Two different things share one name here, and the difference matters.

\textbf{C64 memory to the expansion and back} is the REU's own direct memory
access. There is no command involved: the SDK programs eight registers and the
transfer happens with the processor halted, so by the next instruction it is
already done. Nothing to poll, nothing to wait for.

\textbf{A file to the expansion} is a command, and it is the one that cannot be
done any other way --- the bytes never pass through the C64 at all. Sixteen
megabytes of assets can be loaded without staging them through 64K.

\cmd{ultimate\_reu\_size()} answers how much there is, in 64K banks, and zero
when there is none. It measures rather than asks, because nothing in the
protocol reports it: the expansion aliases, so the first power-of-two boundary
whose write comes back round to the start is the size.

\begin{note}
16MB is the ceiling whatever the machine has, because the expansion's address
registers are 24 bits. A machine with more RAM than that cannot reach the rest
through this interface.
\end{note}

\section{Processor Speed}

Memory-mapped, not a command, and only on an Ultimate~64. The index runs from
0 to 15; 0 is 1MHz and 3 is 4MHz on every machine that has it. Above 3 the same
index means different speeds on different models, so the SDK passes the number
through and does not pretend to know megahertz.

Whether turbo exists at all is the machine owner's decision. Both registers
read \cmd{\$FF} when it has been switched off, which makes availability
testable rather than assumed.

There is a second half to it: bit~7 turns off badlines, which stops the video
chip stealing cycles from the processor. On a fixed loop that is worth about
6.3\% --- measured, and it agrees with the arithmetic.

\section{The Screen Palette}

All sixteen colours, live, on firmware newer than 3.15. This affects the
running palette only --- not flash, not a file --- so a program that crashes
half way through a colour cycle cannot leave the machine looking wrong
permanently.

A whole-palette rotation costs about a quarter of a frame, so colour cycling
directly is comfortably affordable.

\section{Network Sockets}

TCP and UDP, with three things that decide how you use them.

\textbf{A read never waits.} One issued straight after connecting reports
"nothing yet" even when the far end has already replied. Callers poll --- which
is the right answer on a C64, because a blocking read would freeze the machine
for as long as the other end stayed quiet.

\textbf{The end of the stream is \cmd{ULTIMATE\_END}}, the same code a
directory ends with, and it is reported exactly once. After it the handle is
already gone, so do not close it.

\textbf{Connecting can take thirty seconds.} 48 to 75 milliseconds to a host
that is up, and 30.8 seconds to an address with nothing at it before the
firmware gives up.

\begin{listingbox}
for (;;) \{\\
\hspace*{1em}err = ultimate\_net\_read(h, buf, sizeof buf, \&got);\\
\hspace*{1em}if (err == ULTIMATE\_END) break;\hspace{1em}/* far end hung up */\\
\hspace*{1em}if (err != ULTIMATE\_OK)   return err;\\
\hspace*{1em}if (got == 0)             continue;\hspace{1em}/* nothing yet */\\
\hspace*{1em}/* \ldots use buf[0..got-1] \ldots */\\
\}
\end{listingbox}

\section{Fetching Over HTTP}

\cmd{ultimate\_http\_get()} is the whole of the common case: it builds the
request, runs it, and gives the firmware's slot back whether it worked or not.

\begin{listingbox}
uint16\_t got;\\
uint8\_t  err;\\
\\
err = ultimate\_http\_get(url, buf, sizeof buf, \&got);
\end{listingbox}

\cmd{ULTIMATE\_OK} means the server answered below 400. A 404 or a 500 is
\cmd{ULTIMATE\_ERR\_DEVICE} with the number in
\cmd{ultimate\_device\_code()} --- the request itself worked, and the body of
the error page is in your buffer.

\begin{warning}
A URL written as a plain C string literal arrives at the server uppercased.
cc65 translates source characters into PETSCII, and an HTTP path is case
sensitive. Build the URL as bytes, or fold it back at run time.
\end{warning}

\begin{note}
A reply larger than your buffer is truncated and the rest is gone. Nothing in
the protocol offers a range request or a second block, so size the buffer for
what you expect.
\end{note}

% Appendices.

\appendix
\guidechapter{APPENDICES}{%
  \item A --- Every BASIC keyword
  \item B --- Result codes
  \item C --- Where to find the rest
}

\section{Appendix A: Every BASIC Keyword}

Token values are a file format --- a tokenised program stores the number, not
the name --- so this table is append-only and these numbers never change.

\begin{reftable}{llp{0.46\linewidth}}
Token & Keyword & \\
\midrule
\cmd{\$CC} & \bas{uci}    & any command, any target; alone, rescue the interface \\
\cmd{\$CD} & \bas{uerr}   & the last result code \\
\cmd{\$CE} & \bas{udev}   & the raw device code \\
\cmd{\$CF} & \bas{ulen}   & reply length \\
\cmd{\$D0} & \bas{ust\$}  & status text \\
\cmd{\$D1} & \bas{udat\$} & reply as a string \\
\cmd{\$D2} & \bas{ubyte(} & byte n of the reply \\
\cmd{\$D3} & \bas{uw(}    & force 16 bits \\
\cmd{\$D4} & \bas{ul(}    & force 32 bits \\
\cmd{\$D5} & \bas{udos1}  & target 1 \\
\cmd{\$D6} & \bas{udos2}  & target 2 \\
\cmd{\$D7} & \bas{unet}   & target 3 \\
\cmd{\$D8} & \bas{uctrl}  & target 4 \\
\cmd{\$D9} & \bas{uiec}   & target 5 \\
\cmd{\$DA} & \bas{uhttp}  & target 6 \\
\cmd{\$DB} & \bas{uturbo} & set or read the speed index \\
\cmd{\$DC} & \bas{uload}  & load a PRG \\
\cmd{\$DD} & \bas{ubload} & load raw bytes \\
\cmd{\$DE} & \bas{usave}  & memory to a file \\
\cmd{\$DF} & \bas{udir}   & print the directory \\
\cmd{\$E0} & \bas{ustash} & C64 memory into the expansion \\
\cmd{\$E1} & \bas{ufetch} & and back out \\
\cmd{\$E2} & \bas{ureu}   & the expansion's size in banks \\
\end{reftable}

\section{Appendix B: Result Codes}

\begin{reftable}{llp{0.50\linewidth}}
0  & \cmd{ULTIMATE\_OK}               & success \\
1  & \cmd{ULTIMATE\_ERR\_NO\_DEVICE}   & no command interface found \\
2  & \cmd{ULTIMATE\_ERR\_TIMEOUT}     & it did not answer in time \\
3  & \cmd{ULTIMATE\_ERR\_PROTOCOL}    & a transport-level violation \\
4  & \cmd{ULTIMATE\_ERR\_NOT\_SUPPORTED} & this firmware has not got it \\
5  & \cmd{ULTIMATE\_ERR\_INVALID\_ARGUMENT} & your mistake, caught early \\
6  & \cmd{ULTIMATE\_ERR\_IO}          & filesystem or network failure \\
7  & \cmd{ULTIMATE\_ERR\_DEVICE}      & the target reported a failure \\
8  & \cmd{ULTIMATE\_ERR\_TRUNCATED}   & it worked; the buffer was small \\
9  & \cmd{ULTIMATE\_ERR\_ABORTED}     & the command was abandoned \\
10 & \cmd{ULTIMATE\_END}              & no more --- and not a failure \\
\end{reftable}

\cmd{ultimate\_strerror()} turns any of them into a short, stable description.

\section{Appendix C: Where To Find The Rest}

This guide is the friendly half. The reference half lives with the source, and
is generated from the same definitions the code is:

\begin{reftable}{lp{0.55\linewidth}}
\cmd{docs/uci.md} & the protocol itself, with every fact sourced and measured
  behaviour separated from documented behaviour \\
\cmd{docs/asm-abi.md} & the assembly contract, entry point by entry point \\
\cmd{docs/api-design.md} & why the API looks like this \\
\cmd{docs/compatibility.md} & hardware, firmware, speeds, known issues \\
\cmd{docs/generated/} & every constant, every keyword and every command the
  tests exercise --- generated, so never out of step \\
\cmd{bindings/blob/README.md} & the jump table and parameter block, offset by
  offset \\
\end{reftable}


\end{document}
